\documentclass[a4paper]{article}
\usepackage{martimath}
\usepackage{fancyvrb}
\usepackage{longtable,array}
\title{martimath 1.0\\\large Manual d'ús}
\author{Martí Carrasco}
\date{Novembre de 2026}
\begin{document}
\maketitle

\section{Instal·lació i primer document}
Descomprimeix l'arxiu ZIP i deixa els cinc fitxers amb extensió
\texttt{.sty} a la mateixa carpeta del document. La manera més directa
de començar és obrir \texttt{plantilla.tex} i compilar-la amb
\texttt{pdflatex}. Per als índexs, fes dues compilacions.

\begin{Verbatim}[fontsize=\small]
\documentclass[a4paper]{article}
\usepackage{martimath} % 12 pt per defecte
\title{Estructures Algebraiques}
\author{Martí Carrasco}
\begin{document}
\maketitle
\section{Anells}
\begin{definicio}[Ideal]
Sigui $A$ un anell...
\end{definicio}
\end{document}
\end{Verbatim}

La tipografia és \emph{Computer Modern original}, sense substitució.
El text principal té 12 pt per defecte; per canviar-lo, escriu
\texttt{\string\usepackage[fontsize=11pt]\{martimath\}} al preàmbul.
Per defecte hi ha marges d'1,5 cm, seccions i subseccions de vermell fort
(\texttt{\#C00000}), en 14 pt i 12 pt; totes dues conserven les majúscules i
minúscules que escriguis. Les subsubseccions són negres
de 11 pt i paginació a la dreta.

\section{Definicions, teoremes i altres resultats}
Els entorns disponibles són els següents:

\begin{center}
\begin{tabular}{>{\ttfamily}l l}
\hline
definicio & Def: \\
proposicio, prop & Prop: \\
observacio, obs & Obs: \\
teorema, teo & Teorema: \\
lema & Lema: \\
corollari, cor & Cor: \\
exemple & Exemple: \\
comentari & Comentari: \\
demostracio, dem & Dem: i quadrat final \\
\hline
\end{tabular}
\end{center}

El primer claudàtor, si hi és, conté un nom opcional; el segon
controla el disseny: \texttt{inline} (predeterminat),
\texttt{separat}, \texttt{important} i \texttt{index=false}.

\Needspace{15\baselineskip}
\begin{Verbatim}[fontsize=\small]
\begin{prop}
Enunciat sense subtítol.
\end{prop}

\begin{teo}[Nom del teorema][important,separat]
Enunciat amb línia vertical.
\begin{dem}
Demostració dins del teorema.
\end{dem}
\end{teo}
\end{Verbatim}

La demostració pot anar dins o després de l'entorn: en tots dos casos
la barra vertical acaba amb l'enunciat. Una demostració llarga es pot
repartir entre pàgines; els enunciats es protegeixen de talls inadequats.
Les demostracions es mostren en mida petita i blau; les etiquetes dels
resultats i els noms opcionals són vermell fort.

\section{Comandes curtes i opcions globals}
Per a resultats breus, pots escriure \texttt{\string\defi},
\texttt{\string\propi}, \texttt{\string\obsi},
\texttt{\string\teoi}, \texttt{\string\lemai},
\texttt{\string\cori}, \texttt{\string\exi} i
\texttt{\string\comi}. Admeten els mateixos arguments opcionals.

\begin{Verbatim}[fontsize=\small]
\propi[Ideals principals]{Tot ideal principal
    és generat per un sol element.}
\defi{Una definició molt breu.}
\end{Verbatim}

\texttt{\string\def} és una primitiva de TeX: per això no es redefineix.
La comanda curta per a les definicions és \texttt{\string\defi}.

\begin{Verbatim}[fontsize=\small]
\usepackage[fontsize=12pt,mode=normal,
            spacing=normal,colors=true]{martimath}
\usepackage[mode=print]{martimath}
\usepackage[mode=two-up]{martimath}
\end{Verbatim}

\texttt{mode=print} elimina els colors; \texttt{colors=false} fa el
mateix. \texttt{spacing} accepta \texttt{compact}, \texttt{normal}
i \texttt{relaxed}. El valor de \texttt{fontsize} (per defecte 12 pt)
controla la mida del text principal; per exemple, \texttt{10pt} o
\texttt{11pt}, sense canviar la tipografia Computer Modern.

\texttt{mode=two-up} compon dues pàgines A5 verticals de lletra real
12 pt sobre cada full A4 horitzontal, sense reduir-les. Fa servir marges
laterals d'1 cm en cada A5. Per combinar-ho amb el blanc i negre, escriu
\texttt{\string\usepackage[mode=print,twoup=true]\{martimath\}}.

\section{Índexs, figures i fórmules}
Per inserir l'índex general, escriu \texttt{\string\tableofcontents}.

Per inserir el de resultats, escriu \texttt{\string\indexresultats}.

Tots dos índexs són compactes i negres, però tenen enllaços clicables.
No es crea numeració de resultats.
Si un resultat no ha d'entrar a l'índex, usa \texttt{index=false}.

\begin{Verbatim}[fontsize=\small]
\imatge[0.65\linewidth]{dibuix.png}[Llegenda opcional]

\begin{grafica}[xmin=-2,xmax=2,ymin=-2,ymax=4]
\plotfuncion[color=mmblue]{x^2-1}
\plotfuncion[color=black,dashed]{x+0.2}
\end{grafica}
\end{Verbatim}

També tens disponibles \texttt{graphicx}, \texttt{tikz},
\texttt{pgfplots} i \texttt{tikz-cd}. Les fórmules
\verb|\[...\]| no es numeren; si vols una equació numerada,
utilitza l'entorn estàndard \texttt{equation}.

\end{document}
